% cSpell : language = es

\chapter[Tutorial de \LaTeX]{Tutorial de \LaTeX: Fundamentos y Conceptos Clave}
\label{cp:latex-tutorial}

Este capítulo presenta un pequeño manural que introduje al manejo del lenguaje de procesamiento de textos \LaTeX~y los conceptos esenciales para elaborar un TFG. \LaTeX~es una herramienta para crear documentos profesionales a partir de texto plano con comandos de formato, a diferencia de los editores WYSIWYG como Microsoft Word. Estos archivos son procesados por un motor TeX para generar un documento PDF final, lo que te permite centrarte en el contenido mientras \LaTeX~se ocupa del diseño. Aunque este capítulo cubre las características clave, vale la pena aprender \LaTeX~desde el principio. Para una introducción rápida, consulta la serie "Learn LaTeX" de Overleaf \href{https://www.overleaf.com/learn/latex/Learn_LaTeX_in_30_minutes}{Learn LaTeX}.

\section{Citaciones}
\label{sec:citations}

Presentamos dos enfoques distintos para citar entradas en la bibliografía. El primer método corresponde a las citas en el texto, utilizando \verb|\citet{CITA}|, mientras que el segundo emplea \verb|\citep{CITA}| para citas con paréntesis. A continuación hay un ejemplo que muestra ambos usos. Puedes citar múltiples obras en la misma llamada usando \verb|\citep{CITA1, CITA2, ...}|. Para citar sólo el título, usa \verb|\citetitle{CITA}|, y para citar el autor, usa \verb|\citeauthor{CITA}|.

\begin{consejo}
Las citas adecuadas son fundamentales en la escritura académica y aseguran credibilidad, transparencia y avance del conocimiento. Asegúrate de su precisión y pertinencia.
\end{consejo}


\noindent\textbf{Ejemplo:} El principal libro sobre latex es \citet{lamport_latex_1994}, que fue escrito por el creador de \LaTeX~\citeauthor{lamport_latex_1994}, aunque también
es muy interesante el artículo \citep{borbon_latex_2022}.

\section{Referencias cruzadas}

Al igual que con las citas, es recomendable usar referencias en tu documento para citar elementos clave como capítulos, secciones, figuras, tablas o ecuaciones. Para referenciar estos elementos, empieza creando una etiqueta con \verb|\label{ETIQUETA}| y colócala dentro del elemento al que te quieres referir. Una vez creado, puedes usar \verb|\ref{ETIQUETA}| para generar una referencia en el texto. \textbf{Recomendamos encarecidamente usar} \verb|\autoref{LABEL}|, que crea automáticamente un enlace con color según el tipo de elemento. Por ejemplo, una referencia a un capítulo aparecerá como \autoref{cp:introduction} en lugar de Capítulo \ref{cp:introduction}.

\section{Glosario, acrónimos y símbolos}

El documento incluye un glosario, así como listas de acrónimos y símbolos accesibles al inicio. Puedes crear una nueva entrada en \verb|Matter/05-Glossary|, \verb|Matter/| \verb|06-Acronyms| o \verb|Matter/07-Symbols|, según el tipo de entrada. Una vez creada, referencia una entrada de glosario o símbolo con \verb|\gls{ENTRADA}|. Para acrónimos, hay dos formas: \verb|\acrfull{ENTRADA}| para la primera aparición (muestra la definición) y \verb|\acrshort{ENTRADA}| para apariciones posteriores sin repetir la definición.

\vspace{.875em}
\noindent\textbf{Ejemplo:} Utilizar \Gls{latex} para el \acrfull{tfg} es esencial en la \acrfull{eps} de la \acrfull{ceu}. Para facilitar la escritura de la memoria del \acrshort{tfg} la propia \acrshort{eps} facilita esta plantilla.

\section{Figuras}

En \LaTeX, insertar figuras es un proceso sencillo. Para hacerlo, utiliza el entorno \verb|\begin{figure}|. Puedes ajustar el parámetro \verb|width| según tus necesidades, pero asegúrate de seleccionar imágenes de alta calidad (a ser posible en formato \texttt{pdf}, aunque también se admiten en formato \texttt{png} o \texttt{jpg}). También es fundamental proporcionar una leyenda bien redactada. Si procede, incluye citas o referencias que indiquen la fuente de la figura. Usa \verb|\caption{TEXTO}| para la leyenda y \verb|\caption| \verb|[TEXTO_PEQUENO]{TEXTO_GRANDE}| si quieres una versión abreviada para la lista de figuras. Siguiendo estos consejos, podemos crear una figura como la mostrada en \autoref{fig:figure-01}.

\begin{figure}[!htpb]
    \centering
    \includegraphics[width=0.5\linewidth]{Img/Theme/Logotypes/Sello-usp.pdf}
    \caption[Sello del la USP CEU.]{Logotipo de la Universidad San Pablo CEU con forma de sello.}
    \label{fig:figure-01}
\end{figure}

También se pueden insertar subfiguras una a lado de la otra usando 
el entorno \verb|\begin{subfigure}| del paquete \texttt{subcaption}.
Para referirse a las subfiguras, puede usarse \autoref{fig:figure-02.1} y \autoref{fig:figure-02.2}.

\begin{figure}[!htpb]
    \centering
    \begin{subfigure}{0.45\textwidth}
        \centering
        \includegraphics[width=\textwidth]{Img/Theme/Logotypes/Logo-uspceu.pdf}
        \caption{Logotipo azul.}
        \label{fig:figure-02.1}
    \end{subfigure}
    \hspace{.5cm} % Adjust the space as needed.
    \begin{subfigure}{0.45\textwidth}
        \centering
        \includegraphics[width=\textwidth]{Img/Theme/Logotypes/Logo-uspceu-gris.pdf}
        \caption{Logotipo gris.}
        \label{fig:figure-02.2}
    \end{subfigure}
    \caption{Logotipos oficiales de la Universidad San Pablo CEU.}
    \label{fig:figure-02}
\end{figure}

\section{Cuadros}

Los cuadros\footnote{La Real Academia de la Lengua Española recomienda usar el término \emph{cuadro} en lugar de \emph{tabla}.} son esenciales para presentar resultados de forma efectiva. Este capítulo explora técnicas para mostrar información mediante cuadros usando distintos entornos de la plantilla. Definir cuadros en \LaTeX\ puede parecer complejo, pero esta plantilla facilita el proceso.

\begin{consejo}
Los distintos entornos de cuadros deben estar dentro de un entorno \texttt{\textbackslash begin\{table\}} y usar las opciones de flotación \texttt{[!htpb]} para mejorar su colocación. \textbf{Este consejo también aplica al posicionamiento de figuras}.
\end{consejo}

\subsection{Entorno Tabular}

El entorno tradicional \verb|\begin{tabular}| permite crear tablas simples y elegantes. El \autoref{tab:table-01} se genera usando un entorno centrado para mayor énfasis e incorpora la configuración \verb|booktab| para un estilo más sofisticado.

\begin{table}[!htpb]
    \caption{Lenguajes de programación populares.}
    \label{tab:table-01}
    \centering
    \begin{tabular}{llc}
        \toprule
        \textbf{Lenguaje} & \textbf{Extensión} & \textbf{Web} \\ 
        \midrule
        C++          & \texttt{.cpp}   & \url{https://isocpp.org}  \\
        Julia         & \texttt{.jl}    & \url{https://julialang.org}  \\
        Python        & \texttt{.py}    & \url{https://www.python.org}  \\
        R          & \texttt{.r}     & \url{https://www.r-project.org}  \\
        \bottomrule
    \end{tabular}
\end{table}

\subsection{Entorno Tabularx}

Usa el paquete \verb|\begin{tabularx}| para construir una tabla con columnas que se expanden automáticamente. Para ello emplea \verb|\begin{tabularx}{\textwidth}{lX}|, donde \verb|X| será la columna de ancho flexible. Usa \verb|C| para centrarla y \verb|L| y \verb|R| para alinear a la izquierda o derecha, respectivamente. El \autoref{tab:table-02} muestra este entorno en uso.

\begin{table}[!htpb]
    \caption{Una tabla que muestra el uso del entorno \texttt{tabularx}.}
    \label{tab:table-02}
    \begin{tabularx}{\textwidth}{lX}
        \toprule
        \textbf{Header 01} & \textbf{Header 02} \\ 
        \midrule
        Foo Bar Baz & Quisque cursus, metus vitae pharetra auctor, sem massa mattis sem, at interdum magna augue eget diam. \\
        Ipsum Dolor & Vestibulum ante ipsum primis in faucibus orci luctus et ultrices posuere cubilia Curae; Curabitur aliquet quam id dui. \\
        Dolor Sit & Phasellus condimentum elementum justo, quis interdum est sagittis ac. Vestibulum non arcu sit amet justo lobortis semper. \\
        Amet Consectetuer & Integer nec odio praesent libero sed cursus ante dapibus diam sed nisi vestibulum non arcu. \\
        % Consectetuer Adipiscing & Nulla quis sem at nibh elementum imperdiet. Duis sagittis ipsum. Praesent mauris. \\
        \bottomrule
    \end{tabularx}
\end{table}

\subsection{Entorno Longtable}

Cuando trabajes con tablas muy largas puede ser necesario dividirlas en varias páginas. En \LaTeX\ se logra con el entorno \verb|\begin{longtable}|. Este entorno es algo más complejo porque requiere definir la cabecera dos veces: una para la primera aparición y otra para las páginas siguientes, de modo que el lector identifique correctamente las columnas. Consulta el \autoref{tab:table-03} para un ejemplo detallado.

\begin{longtable}[c]{lcc}
\caption{Asignaturas del grado en Ingeniería Matemática.}
\label{tab:table-03} \\
\toprule
\textbf{Asignatura} & \textbf{Curso} & \textbf{Créditos} \\ \midrule
\endfirsthead
%
\multicolumn{3}{c}%
{{\textit{\bfseries Cuadro \thetable\ (continúa de la página anterior).}}} \\
\toprule
\textbf{Asignatura} & \textbf{Curso} & \textbf{Créditos} \\ \midrule
\endhead
%
\bottomrule
\addlinespace[1mm]
\multicolumn{3}{r}%
{{\textit{Continúa en la siguiente página.}}} \\
\endfoot
\bottomrule
%
\endlastfoot
%
%
% --- Primer curso - 1º semestre ---
Fundamentos de Matemáticas & 1º & 3 \\
Análisis I & 1º & 6 \\
Álgebra Lineal & 1º & 9 \\
Matemática Discreta & 1º & 6 \\
Programación & 1º & 6 \\

% --- Primer curso - 2º semestre ---
Análisis II & 1º & 6 \\
Probabilidad y Estadística & 1º & 6 \\
Algoritmos y Estructuras de Datos & 1º & 9 \\
Claves de Historia Contemporánea & 1º & 6 \\
Proyecto I & 1º & 3 \\

% --- Segundo curso - 1º semestre ---
Análisis III & 2º & 6 \\
Ecuaciones Diferenciales y en Diferencias & 2º & 6 \\
Bases de Datos & 2º & 6 \\
Estadística Inferencial & 2º & 6 \\
Fundamentos Económicos & 2º & 6 \\

% --- Segundo curso - 2º semestre ---
Geometría Diferencial & 2º & 3 \\
Métodos Numéricos I & 2º & 6 \\
Electrónica Digital y Arquitectura de Ordenadores & 2º & 6 \\
Ecuaciones en Derivadas Parciales & 2º & 6 \\
Análisis de Datos & 2º & 6 \\
Proyecto II & 2º & 3 \\

% --- Tercer curso - 1º semestre ---
Sistemas Operativos y Redes de Ordenadores & 3º & 6 \\
Optimización & 3º & 6 \\
Métodos Numéricos II & 3º & 6 \\
Hombre y Mundo Moderno & 3º & 6 \\

% --- Optativas (3º) ---
Sistemas Dinámicos (Optativa) & 3º & 3 \\
Análisis Funcional (Optativa) & 3º & 3 \\
Topología (Optativa) & 3º & 3 \\
Variable Compleja y Análisis de Fourier (Optativa) & 3º & 6 \\

% --- Tercer curso - Mención en Inteligencia Artificial (segundo semestre) ---
Aprendizaje Automático & 3º & 6 \\
Computación en paralelo & 3º & 6 \\
Lógica Formal & 3º & 6 \\
Teoría de la Computación & 3º & 6 \\
Fundamentos de la Inteligencia Artificial & 3º & 6 \\

% --- Tercer curso - Mención en Análisis Cuantitativo y Finanzas (segundo semestre) ---
Aprendizaje Automático & 3º & 6 \\
Computación en paralelo & 3º & 6 \\
Cálculo Estocástico & 3º & 6 \\
Matemática Financiera I & 3º & 6 \\
Matemáticas Actuariales & 3º & 6 \\

% --- Cuarto curso (Mención en Inteligencia Artificial) - 1º semestre ---
Programación Lógica & 4º & 6 \\
Programación Funcional & 4º & 6 \\
Percepción Computacional & 4º & 6 \\
Procesamiento de Lenguaje Natural & 4º & 6 \\
Administración de sistemas & 4º & 6 \\

% --- Cuarto curso (Mención en Inteligencia Artificial) - 2º semestre ---
Doctrina Social de la Iglesia & 4º & 6 \\
Prácticas externas & 4º & 12 \\
Trabajo Fin de Grado & 4º & 9 \\

% --- Optativas 4º (Mención IA) ---
Computación Cuántica (Optativa) & 4º & 3 \\
Aprendizaje profundo (Optativa) & 4º & 3 \\
Teoría de la señal (Optativa) & 4º & 3 \\
Ingeniería del Software (Optativa) & 4º & 3 \\
Grandes libros (Optativa) & 4º & 3 \\
Ética y Deontología (Optativa) & 4º & 3 \\

% --- Cuarto curso (Mención en Análisis Cuantitativo y Finanzas) - 1º semestre ---
Matemática Financiera II & 4º & 6 \\
Minería de Datos y Big Data & 4º & 6 \\
Modelos de Riesgo Cuantitativo & 4º & 6 \\
Teoría y optimización de carteras & 4º & 6 \\
Series Temporales & 4º & 6 \\

% --- Optativas 4º (Mención ACyF) ---
Computación Cuántica (Optativa) & 4º & 3 \\
Aprendizaje profundo (Optativa) & 4º & 3 \\
Procesos Estocásticos (Optativa) & 4º & 3 \\
Criptografía y Blockchain (Optativa) & 4º & 3 \\
Grandes libros (Optativa) & 4º & 3 \\
Ética y Deontología (Optativa) & 4º & 3 \\

\end{longtable}

\subsection{Tablas complejas}

Crear tablas complejas en \LaTeX\ puede resultar complicado. Recomendamos usar \href{https://www.tablesgenerator.com/}{Table Generator}, con el que diseñas la tabla y copias el resultado al documento. Así se simplifica el proceso y se evita errores. Ten en cuenta que la tabla debe ser comprensible y no excesivamente compleja, pues \textbf{la complejidad puede dificultar la comprensión}. Por ejemplo, el \autoref{tab:table-04} muestra una tabla con detalles intrincados.

\begin{table}[!htpb]
    \caption{Especificaciones de hardware para componentes clave de un ordenador.}
    \label{tab:table-04}
    \centering
    \begin{tabular}{lcc}
        \toprule
        \multirow{2}{*}{\textbf{Component}} & \multicolumn{2}{c}{\textbf{Specifications}} \\
        \cmidrule(lr){2-3}
        & \textbf{Characteristic} & \textbf{Supported} \\
        \midrule
        \multirow{4}{*}{CPU} & Core Count (e.g., 8 Cores) & $\checkmark$ \\
        & Clock Speed (e.g., 3.6 GHz) & $\checkmark$ \\
        & Hyper-Threading & $\checkmark$ \\
        & Integrated Graphics & - \\
        \midrule
        \multirow{4}{*}{GPU} & CUDA Cores (e.g., 5120) & $\checkmark$ \\
        & Base Clock (e.g., 1.5 GHz) & $\checkmark$ \\
        & Ray Tracing Support & $\checkmark$ \\
        & Multi-GPU Support (SLI/CrossFire) & - \\
        \midrule
        \multirow{4}{*}{Memory} & Type (e.g., DDR5, GDDR6) & $\checkmark$ \\
        & Capacity (e.g., 16 GB) & $\checkmark$ \\
        & Memory Bandwidth (e.g., 448 GB/s) & $\checkmark$ \\
        & ECC Support & - \\
        \midrule
        \multirow{3}{*}{Motherboard Features} & PCIe 5.0 Support & $\checkmark$ \\
        & Wi-Fi 6E & $\checkmark$ \\
        & Thunderbolt 4 & - \\
        \bottomrule
    \end{tabular}
\end{table}

\section{Listas}

Crear listas en \LaTeX\ es sencillo y ofrece varias opciones según lo que necesites. Usa \verb|\begin{itemize}| para listas con viñetas o \verb|\begin{enumerate}| para listas numeradas. A continuación se muestra un ejemplo con \verb|\begin{itemize}|.


\begin{itemize}
  \item Las entradas de la lista comienzan con el comando \verb|\item|.
  \item Las entradas individuales se indican con un punto negro, conocido como viñeta.
  \item El texto en las entradas puede tener cualquier longitud.
\end{itemize}

Como se ha mencionado anteriormente, puedes generar una lista numerada utilizando el entorno \verb|\begin{enumerate}|. Aquí tienes un ejemplo:

\begin{enumerate}
  \item Los elementos se numeran automáticamente.
  \item Los números comienzan en 1 con cada uso del entorno \verb|enumerate|.
  \item Otra entrada en la lista.
\end{enumerate}

También puedes anidar entradas de lista creando una lista dentro de otra lista del mismo tipo. Aquí tienes un ejemplo:

\begin{enumerate}
    \item Primer elemento de nivel
    \item Primer elemento de nivel
    \begin{enumerate}
        \item Segundo elemento de nivel
        \item Segundo elemento de nivel
    \begin{enumerate}
        \item Tercer elemento de nivel
        \item Tercer elemento de nivel
    \end{enumerate}
    \end{enumerate}
\end{enumerate}

\begin{consejo}
Ten en cuenta que las etiquetas cambian automáticamente incluso si usas el mismo tipo de entorno para distintas listas. \textbf{No es necesario preocuparse por cambiar el entorno para obtener un formato distinto.}
\end{consejo}

También puedes modificar la etiqueta de tu lista a algo completamente diferente que se ajuste a tus necesidades. Para lograr esto, inserta un nuevo \verb|\item| y encierra la etiqueta deseada entre corchetes. Por ejemplo, \verb|\item[!]| resultará en un signo de exclamación como tu nueva etiqueta. A continuación se muestran algunos ejemplos de etiquetas modificadas.

\begin{itemize}
  \item Este es mi primer punto
  \item Otro punto que quiero hacer
  \item[!] ¡Un punto para exclamación!
  \item[$\blacksquare$] Hacer el punto justo y cuadrado.
  \item[] ¿Una etiqueta en blanco?
\end{itemize}

Finalmente, puedes crear una lista de descripción. A diferencia de tener un punto de viñeta o una etiqueta numerada, una lista de descripción te permite usar descripciones personalizadas que se ajusten a tu lista. En el ejemplo a continuación, hay tres entradas \verb|\item|: una sin etiqueta y dos con descripciones.

\begin{description}
    \item[Item 1:] Este es el primer elemento con una descripción.
    \item[Item 2:] Otro elemento con una descripción diferente.
    \item Un elemento sin una etiqueta específica.
\end{description}

\section{Listados de código}

En ocasiones, es necesario incluir código fuente en el documento. Para ello puedes usar dos entornos anidados: \verb|\begin{listing}| para el entorno con título y etiqueta, y \verb|\begin{minted}| para el resaltado de sintaxis. El \autoref{listing:c-code} muestra un ejemplo en C.

\begin{listing}[!htpb]
\caption{Hello world in C.}
\label{listing:c-code}
\begin{minted}{c}
#include <stdio.h>
int main() {
   printf("Hello, World!"); /* printf() outputs the quoted string */
   return 0;
}
\end{minted}
\end{listing}

El código anterior fue insertado directamente en el documento. Como alternativa, puedes cargarlo desde un archivo externo con \verb|\inputminted{LENGUAJE}{ARCHIVO}|, colocando ese comando dentro de un entorno \verb|\begin{listing}|. El \autoref{listing:julia-code} muestra un ejemplo usando Julia desde el fichero \texttt{Code/Factorial.jl}.

\begin{listing}[!htpb]
\caption{Factorial in Julia.}
\label{listing:julia-code}
\inputminted{julia}{Code/Factorial.jl}
\end{listing}

En algunos casos, cuando simplemente deseas resaltar un comando específico, se recomienda no usar \verb|listing| o \verb|minted|. En su lugar, debes utilizar el comando \verb|\verb| para resaltado en línea o el entorno \verb|\begin{verbatim}| para secciones más largas de código resaltado. A continuación se proporciona un ejemplo de una sección extensa de \verb|verbatim|, que demuestra cómo crear un \verb|listing| con un código de entrada:

\begin{verbatim}
\begin{listing}[!htpb]
    \inputminted{LENGUAJE}{ARCHIVO}
    \caption{TEXTO}
    \label{ETIQUETA}
\end{listing}
\end{verbatim}

A veces es necesario mostrar código más extenso que ocupa más de una página. Para ello, utiliza el entorno \verb|\begin{longlisting}|. Este entorno dividirá fácilmente tu código en varias páginas para una mejor legibilidad, sin que tengas que preocuparte por el tamaño de tu código. A continuación se muestra un ejemplo en el \autoref{listing:python-code}.

\begin{longlisting}
\caption{Ejemplo de funciones en Python.}
\label{listing:python-code}
\begin{minted}{python}
from typing import List

def factorial(n: int) -> int:
    """Calculate the factorial of a number (recursive)."""
    if n < 0:
        raise ValueError("factorial is not defined for negative integers")
    return 1 if n <= 1 else n * factorial(n - 1)

def fibonacci(n: int) -> int:
    """Calculate the nth Fibonacci number (recursive)."""
    if n == 0:
        return 0
    if n == 1:
        return 1
    return fibonacci(n - 1) + fibonacci(n - 2)

def gcd(a: int, b: int) -> int:
    """Greatest common divisor using Euclid's algorithm."""
    a, b = abs(a), abs(b)
    while b:
        a, b = b, a % b
    return a

def primes_up_to(limit: int) -> List[int]:
    """Return a list of all prime numbers up to limit (inclusive)."""
    if limit < 2:
        return []
    primes: List[int] = []
    for i in range(2, limit + 1):
        is_prime = True
        for p in primes:
            if i % p == 0:
                is_prime = False
                break
            if p * p > i:
                break
        if is_prime:
            primes.append(i)
    return primes

def example_function(x: int):
    """Example function to mirror the Lisp version's behavior."""
    result = [factorial(x), fibonacci(x), gcd(x, 10), primes_up_to(x)]
    print(f"Factorial of {x}: {result[0]}")
    print(f"Fibonacci of {x}: {result[1]}")
    print(f"GCD of {x} and 10: {result[2]}")
    print(f"Primes up to {x}: {result[3]}")
    return result


if __name__ == '__main__':
    example_function(10)
\end{minted}
\end{longlisting}

\section{Ecuaciones}

Al escribir ecuaciones y expresiones matemáticas, \LaTeX~es una herramienta potente y versátil. Puedes insertar una fórmula en línea con \verb|\(FORMULA\)| o usar \verb|\begin{equation}| para mostrarla en modo matemático con numeración. Si no deseas numerarla, usa \verb|\[FORMULA\]|.

\vspace{.875em}
\textbf{Ejemplo:} In physics, the mass-energy equivalence is expressed by the equation \(E=mc^2\), discovered in 1905 by Albert Einstein. In natural units ($c = 1$), the formula (\ref{eq:equation-01}) expresses the identity:

\begin{equation}
\label{eq:equation-01}
E=m
\end{equation}

\textbf{Ejemplo:} A continuación se muestra una ecuación -- \textit{sin numeración} -- para la función de pérdida regularizada en el aprendizaje supervisado, combinando la pérdida de predicción promedio sobre el conjunto de datos de entrenamiento y un término de regularización $L_2$ para prevenir el sobreajuste:

\[
\mathcal{L}(\boldsymbol{\theta}) = \frac{1}{N} \sum_{i=1}^{N} \ell(y_i, f(\mathbf{x}_i; \boldsymbol{\theta})) + \lambda \|\boldsymbol{\theta}\|_2^2
\]

La escritura de ecuaciones puede ser un poco desafiante, por lo que recomendamos utilizar un editor en línea, como el \href{https://latexeditor.lagrida.com/}{LaTeX Equation Editor}. Simplemente construye tus fórmulas allí y cópialas y pégalas en tu documento, ya sea en línea o en un bloque matemático, como se muestra arriba.

\section{Notas a pie de página}

A veces es importante presentar información no central en una nota a pie de página. En \LaTeX\ puedes hacerlo con \verb|\footnote{TEXTO}|; el texto aparecerá en la parte inferior de la página\footnote{Esta es una nota a pie de página simple.}.

Si deseas utilizar notas a pie de página dentro de tablas, es mejor reconsiderarlo, ya que \LaTeX\ no proporciona una forma fácil de manejarlas. En su lugar, puedes colocar un ``*'' donde quieras que aparezca la referencia de la nota a pie de página. Luego, debajo de la tabla \textbf{pero antes de finalizar el entorno de la tabla}, coloca el ``*'' junto con el texto de la nota a pie de página. Esto creará una nota a pie de página similar, pero aparecerá debajo de la tabla en lugar de en la parte inferior de la página.